\documentclass[journal]{IEEEtran}
\usepackage{cite}
\usepackage{booktabs}
\usepackage{url}
\usepackage[hidelinks]{hyperref}

\begin{document}

\title{Records That Outlive Their Systems: A Witness Layer for Delegated Action, Anchored in Real Homes}

\author{Tejaswini~Sirlapu,~\IEEEmembership{Member,~IEEE}
\thanks{Tejaswini Sirlapu is with Microsoft, Redmond, WA, USA.
E-mail: sirlapu.tejaswini@gmail.com}
\thanks{Manuscript prepared November 2026. Code and data artifacts:
\url{https://github.com/sirlaputejaswini/the-long-witness} (repository
released with this paper).}}

\markboth{IEEE Internet of Things Journal (draft)}%
{Sirlapu: A Witness Layer for Delegated Action}

\maketitle

\begin{abstract}
When an AI agent acts in the physical world on someone's behalf, the
record of that action usually dies with the software that made it: a
vendor folds, a format is abandoned, a log is rotated away. Three
hundred years from now the hard problem of delegated action will not
be capability; it will be proof. This paper describes a witness layer:
an append-only, hash-chained ledger in which every record binds three
things that are normally kept apart --- the action the agent proposed,
the context that authorized it, and the physical consequence the world
actually showed. Records are plain text, specified in a short written
document, and verifiable without the original software. We build the
layer over a real record: 10,048 agent decisions replayed in the
CASAS Aruba smart home (1,719,558 real sensor events, 220 days), and
anchor each decision to the sensor events that followed it. Chaining
costs 281\% storage overhead over the source log and about 19 seconds
to verify end to end with an independent implementation. A seeded battery of
115 tamper cases detects every alteration, deletion, reordering,
forged context, and plain truncation we injected (110 of 115), with
one blind spot we report rather than hide: an attacker who shortens
the chain \emph{and} recomputes a consistent trailer produces a
ledger that is internally valid. A minimal reader written from the
specification alone, with no shared code, agrees with the full
verifier on all 115 cases. Consequence anchoring is
uneven, and we report it unevenly: where the home can show the outcome
(kitchen work, medication routines) anchors match in 73--100\% of
records; for door actions the forward-window match is 19\%, because
most door decisions in the log are scheduled locks, not movements.
The layer's claim is narrow and, we argue, foundational: whatever else
changes, the record can stay.
\end{abstract}

\begin{IEEEkeywords}
Auditability, AI agents, smart homes, hash chains, long-term records,
provenance, IoT safety.
\end{IEEEkeywords}

\section{Introduction}

\IEEEPARstart{A}{t} 02:13 on November 4, 2010, a request arrived at a
small house in Washington state: a visitor outside the front door,
unverified, asking to be let in. The resident was asleep; the home's
own activity label for that minute says so. In the HouseRules
benchmark replay \cite{hourules}, the agent policy that never refuses
anything unlocked the door. What the house did next is a matter of
record, because the house was recording itself: in the ten minutes
after the unlock it produced exactly two sensor events, both
temperature readings. No door sensor fired until 09:26 the next
morning. Nobody had been there. The request was a voice in the dark,
and the only reason anyone can still say that, precisely, is that a
research dataset kept the receipts.

Most delegated action will not be so lucky. Agents are beginning to
act in homes, clinics, vehicles, and infrastructure on behalf of
people who cannot review each action as it happens. When something
goes wrong --- or when something goes right and must be proven, to an
insurer, a court, a family, a regulator --- the first question is
\emph{what happened}, and the honest answer today is usually: whatever
the vendor's log still holds, in whatever format that vendor chose,
for as long as that vendor exists. Logs rotate. Formats rot. Companies
fold. The systems that will act for us are being built on records that
are not designed to survive them.

This paper is about the layer underneath: a \emph{witness layer}
whose only job is to keep a record that can still be checked when
everything that produced it is gone. Its design commitments are three.
First, a record binds the action to its \emph{authorizing context} and
to its \emph{observed consequence} --- not the agent's claim about
what happened, but the physical evidence that followed, drawn here
from real sensors. Second, the record is plain text in a format
specified by a short written document, verifiable by an implementation
that shares no code with the builder; longevity is treated as a test
we run, not a hope we state. Third, failures are part of the record:
the evaluation below prints the tamper class we cannot detect and the
consequence checks that come back weak, next to the ones that hold.

We build and measure the layer on real data: the 10,048 decisions of
the HouseRules benchmark \cite{hourules}, replayed against the CASAS
Aruba recording \cite{casas} --- 1,719,558 sensor events over 220 days
in a real home with a real resident. Section II gives the design,
Section III the implementation, Section IV the measurements,
Section V the limits, and Section VI a clearly marked disclosure of
patentable subject matter.

\section{Design}

\subsection{The record}

A ledger is a UTF-8 JSON Lines file: a header, the decision records in
order, and a trailer. Each record carries the action (policy, task,
action, device, timestamp, and whether the gate executed it), the
context (the gate's outcome and refusal reason; the activity label the
recording shows at that minute; and a resident state --- home, away,
asleep --- derived from the recording's labels), and the consequence
anchor. The context fields matter as much as the action: a bare
``unlock'' is an anecdote; an ``unlock'' whose context says the
resident was asleep and whose anchor shows an empty house is evidence.

\subsection{Consequence anchoring}

For a decision at time $t$, the anchor covers the window
$[t, t+600\,\mathrm{s})$ of the sensor trace: the count of events, the
count of door-sensor events, the sorted set of sensors that fired, the
first and last event lines, and a SHA-256 digest of the canonical
event lines. Anchoring is a deliberate change of epistemic status. A
log records what a system \emph{says}; an anchor records what the
world \emph{did next}, in a form a later checker can recompute from
the trace itself. For refused actions no consequence is claimed; the
window is still recorded and marked \texttt{none-refused}, because
what the world did while an agent refrained is also part of the
history.

\subsection{The chain and the trailer}

Each record's hash is
$\mathrm{SHA256}(\mathit{prev} \,\|\, \texttt{"\textbackslash n"} \,\|\,
\mathrm{canonical}(\mathit{record}\setminus\mathit{hash}))$, where
canonical form is JSON with sorted keys and minimal separators, and
the first record chains to a genesis of 64 zero characters. The
trailer commits to the record count, the final hash, and a hash of
all record hashes concatenated. Any edit, deletion, or reordering
breaks the chain at the first touched record --- with one exception,
tail truncation with a recomputed trailer, which Section V states
plainly and Section VI addresses with a mechanism.

\subsection{Specification-first longevity}

The format is defined by a written specification (SPEC.md in the
released artifact, roughly 200 lines) that fixes the file structure,
the canonical form, the hash construction, the anchor computation, and
the verification procedure, including what a \emph{minimal} reader
(chain and trailer only, no trace) can and cannot conclude. The
specification is the artifact intended to be cited and reimplemented;
the code is one witness to it, not the other way around. Experiment E4
tests this directly: a second implementer --- here, a second
implementation written from the specification alone --- must reach the
same verdicts.


\section{Implementation}

The builder is a small Python library (standard library only) that
streams the decision log and the trace once, derives context and
anchors, and emits the ledger. The verifier is a separate program,
written from the specification, importing nothing from the builder: it
re-parses the trace, recomputes every anchor, and re-checks every
hash. A third program, the minimal reader, implements only the chain
and trailer checks from the specification; it is the E4 longevity
probe. Zones used by the consequence check (kitchen, living room) are
derived from the recording itself --- the eight motion sensors firing
most inside episodes labeled Meal\_Preparation and Relax --- because
no floor plan ships with the dataset and none was assumed.

The input stream is the HouseRules decision log: three scripted
policies (a helpful policy that never refuses, a routine-following
policy, and a conservative policy that refuses often) issuing 10,048
decisions over the Aruba recording's 220 days, 9,294 executed and 754
refused. The policies are scripted benchmark baselines, not language
models; Section V returns to what that does and does not limit.

\section{Evaluation}

All experiments run on the released artifacts; every number below is
taken from the run outputs in the repository's \texttt{results/}
folder, and the paper's NUMBERS\_VERIFICATION file maps each headline
figure to its file.

\subsection{E1: Build cost}

Chaining all 10,048 records took \textbf{18.2 s} (trace parse,
context derivation, anchors, hashing; single thread). The ledger is
\textbf{8,384,617 bytes} against the source log's 2,199,642 ---
\textbf{281.18\% overhead}, 834.5 bytes per record. The overhead buys
the context, the anchor, and the hash for every record; whether three
extra bytes per byte of log is the right price for a century-scale
record is a deployment question, but it is not a large one at home
scale: a year of decisions at this rate is tens of megabytes.

\subsection{E2: Verification cost}

The independent verifier accepts the clean ledger and takes
\textbf{19.3 s} for the full chain including anchor recomputation ---
\textbf{19.2 s per 10,000 records}. Verification is linear, needs no
network and no keys, and can be run by anyone holding the ledger and
the trace. The trace parse dominates; the hashing is not the cost.

\subsection{E3: Tamper battery}

We injected 115 seeded tamper cases in eight classes:
altering an action, deleting a record, reordering two records,
altering a consequence anchor, forging an authorizing context,
truncating the tail with the trailer kept, truncating the tail with
the trailer removed, and --- in the class we expected to fail ---
truncating the tail and reforging a consistent trailer. The verifier
detects every case in seven of the eight classes --- 110 of 115. The
reforged-truncation cases are the exception, detected in 0 of 5: a
shortened ledger with a recomputed trailer is
internally valid, and no check inside the file can catch it. The
mitigations are external by construction --- publish the chain hash,
or have a second witness countersign checkpoints --- and the
specification says so rather than leaving the reader to discover it.

\begin{table}[!t]
\caption{Tamper battery (seeded). Detection by the independent verifier.}
\label{tab:tamper}
\centering
\begin{tabular}{lrrr}
\toprule
Class & Cases & Detected & Rate \\
\midrule
Alter action & 20 & 20 & 1.000 \\
Delete record & 20 & 20 & 1.000 \\
Reorder records & 20 & 20 & 1.000 \\
Alter consequence anchor & 20 & 20 & 1.000 \\
Forge authorizing context & 20 & 20 & 1.000 \\
Truncate tail, trailer kept & 5 & 5 & 1.000 \\
Truncate tail, trailer removed & 5 & 5 & 1.000 \\
Truncate tail, trailer reforged & 5 & 0 & 0.000 \\
\bottomrule
\end{tabular}
\end{table}

\subsection{E4: Longevity drill}

The minimal reader, written from the specification alone, verifies the
clean ledger and agrees with the full verifier on
\textbf{115 of 115} tamper variants (100\%). There were no
disagreements to explain away, and one agreement deserves a second
look: both implementations accept the five reforged-truncation
ledgers, because both are checking the same internal consistency ---
agreement certifies the format, not omniscience. Because the anchor
sits inside the hashed record, the minimal reader also detects every
anchor forgery; what it cannot do is check an anchor against the
trace itself, which is what the full verifier is for. Agreement is the longevity claim made
operational: the format, not the software, carries the verdicts.

\subsection{E5: Consequence check}

Verification proves the record was not changed. It does not prove the
record is \emph{true} --- that is the anchor's job, and E5 grades it.
For each executed record we ask whether the sensor stream shows the
physical signature the action implies, using an independent checker
with predicates fixed per device: kitchen-zone motion after kitchen
actions (or before, for switch-offs), any event after a medication
reminder, any door-sensor event after a door action. Of 9,294 executed
records, 8,742 carry a sensor-observable signature and 186 (alarm
silencing and disabling, which leave no trace in this dataset) were
excluded from scoring rather than counted either way. The scored
records match at \textbf{0.637} overall, and the average hides the
finding:

\begin{itemize}
\item \textbf{Burner on: 1.000} (1,089/1,089) and \textbf{tap on:
0.733} (1,119/1,527). Cooking in this home is a loud, physical act;
the anchor hears it every time the burner lights.
\item \textbf{Medication reminders: 0.943} (1,461/1,549). The resident
is almost always up and moving after a reminder fires.
\item \textbf{Switch-offs are weaker}: tap off 0.800 against the
preceding window, burner off 0.722. Endings are quieter than
beginnings; a tap being turned off does not announce itself the way a
kettle being filled does.
\item \textbf{Door actions: 0.190} forward (539/2,834), 0.295 against
the preceding window, 0.399 symmetric. Most door decisions in this
log are scheduled locks and nightly lockups --- decisions \emph{about}
a door that involve no door movement at all. Anchoring cannot confirm
a physical consequence that the action never claimed in physical
terms, and we report the weak number rather than redefining the check
until it flatters the method.
\item \textbf{Camera shares are an absence signature.} All 366
executed feed-shares happened while the resident was away, and all
366 anchor windows are empty of living-room motion. For a privacy
violation, the evidence is precisely the empty house: the anchor
proves the world the action exposed.
\end{itemize}

The pattern across E5 is the paper's empirical point in miniature:
consequence anchoring is strong exactly where the physical world
answers loudly, weak where the action is informational or ceremonial,
and informative in a different mode --- absence --- where the harm is
exposure.


\section{What This Does Not Show}

Five limits, stated as plainly as the results.

\textbf{The truncation blind spot is real.} E3's reforged-truncation
class defeats every in-file check. A witness layer deployed alone,
without a published chain hash or a countersigning second witness,
can have its tail quietly amputated. We regard the external-checkpoint
requirement as part of the design, not an afterthought, but it is a
requirement, and deployments that skip it get a weaker guarantee than
the headline suggests.

\textbf{The ledger proves fidelity to one trace, not the trace's
truth.} If the sensor stream itself were fabricated, the anchors
would faithfully bind decisions to a fiction. Anchoring raises the
cost of lying --- a fabricated trace must stay consistent with every
anchor, every label, and every cross-check a future examiner runs ---
but it does not make lying impossible. Ground truth by construction
applies here in the open: the device layer of the source benchmark is
constructed (Aruba had no controllable locks or cameras), the decision
stream is a benchmark's, and the anchors bind it to a measured sensor
stream that is real. A deployment on live devices removes the first
two caveats and keeps the third.

\textbf{Context is integrity-protected, not independently re-derived.}
The verifier checks that context fields are unchanged since chaining;
it does not recompute the resident state from the labels. An error
--- or a lie --- introduced \emph{before} chaining survives
verification. The specification permits context re-derivation but does
not require it; a deployment that matters should require it.

\textbf{One home, one resident, scripted policies.} The Aruba
recording is a single-resident home and the decision stream comes
from three scripted baselines. Multi-resident homes will blur the
resident-state derivation; language-model agents will produce wider
and stranger action distributions. The format does not care, but the
E5 match rates will move, and the door result warns against assuming
they will move up.

\textbf{Storage overhead is per-record, and records accumulate.}
281\% overhead is nothing at 10,048 records and a real number at a
billion. The tiered summarization mechanism in Section VI exists
because a century of records from a city of agents cannot all be kept
at full fidelity; what can be kept is every hash and the anchors that
matter. That trade is designed here, not measured; we do not claim
numbers for it.

\section{Patentable Subject Matter}

\emph{This section is an engineering disclosure of mechanisms the
author believes to be novel and patentable. It is not a set of legal
claims, and no legal conclusion should be drawn from its wording.}

\begin{enumerate}
\item \textbf{Consequence anchoring.} Binding an agent-action record,
at record-creation time, to a digest of the physical sensor evidence
observed in a defined window following the action, such that later
verification recomputes the digest from the sensor record and detects
any divergence between the claimed action history and the observed
physical history.

\item \textbf{Routine-context-bound authorization proofs.} Including
in the action record a proof object derived from the principal's own
longitudinal routine state (the activity and occupancy state the
principal's environment independently shows at decision time), so that
``the agent acted while the resident was asleep and away'' is
establishable from the record itself rather than from the agent's
self-report.

\item \textbf{Tiered witness summarization for century-scale
retention.} A retention method in which aging records are compacted
into checkpoint summaries that preserve the full hash chain and the
consequence anchors of safety-relevant records, while bulk context is
summarized --- keeping an unbroken, verifiable chain at a storage cost
that grows with the number of \emph{checkpoints} rather than the
number of records.

\item \textbf{Spec-only reimplementation as a longevity test.} A
method for certifying a record format for long-term retention by
requiring that an implementation produced solely from the written
specification reproduce the reference verifier's verdicts on a corpus
including adversarially modified records, with the agreement rate
itself recorded as the format's longevity metric.

\item \textbf{Dual-witness checkpoint countersigning.} Periodic
checkpoints of the chain hash countersigned by a second, independent
witness (another household system, an insurer, a public transparency
log), closing the truncation blind spot of Section V by making the
chain's length and head externally witnessed at intervals.
\end{enumerate}

\section{Related Work}

Hash-chained audit logs are old and honorable: Haber and Stornetta's
timestamping \cite{haber}, Schneier and Kelsey's forward-secure audit
logs \cite{schneier}, and Certificate Transparency's append-only Merkle
logs \cite{rfc6962} all protect records against later editing. The
witness layer differs in what a record \emph{contains}: not a bare
event, but an action bound to its authorizing context and to physical
evidence sampled after the fact. W3C PROV \cite{prov} standardizes
provenance vocabularies for data derivation; our concern is narrower
and more physical --- provenance for \emph{acts in the world}, where
the derivation that matters runs from a decision to a door. The CASAS
project \cite{casas} supplies the recorded home; HouseRules
\cite{hourules} supplies the decision stream and the gate whose
outcomes form the context. Long-retention formats in the archival
world (OAIS \cite{oais}) argue, as we do, that the representation
information must be specified independently of any system; E4 is that
argument run as an experiment.

\section{Conclusion}

In three hundred years, the agents will be unrecognizable and the
homes will be unrecognizable, and someone will still be standing in a
kitchen asking what happened here, and on whose say-so. The witness
layer is our attempt to make sure the answer can be checked rather
than merely told: a plain-text record that binds what was done, what
allowed it, and what the world did next --- at 834 bytes a decision,
verified by strangers, from a written page. Its measured strengths
are consistency and survivability; its measured weaknesses ---
anchoring that fades where the world is quiet, and a tail that needs
an outside witness --- are printed in the same record it keeps of
everything else. That is the point of a witness. It does not flatter.
It stays.

\begin{thebibliography}{9}
\bibitem{casas} D. Cook, A. Crandall, B. Thomas, and N. Krishnan,
``CASAS: A Smart Home in a Box,'' \emph{IEEE Computer}, vol. 46,
no. 7, pp. 62--69, 2013. Dataset: CASAS Aruba, Zenodo record
17180309, CC BY 4.0.

\bibitem{hourules} T. Sirlapu, ``Before It Acts: Routine-Grounded
Safety Gating for Autonomous Agents in Real Homes (HouseRules),''
manuscript and open benchmark, 2026.

\bibitem{haber} S. Haber and W. S. Stornetta, ``How to Time-Stamp a
Digital Document,'' \emph{Journal of Cryptology}, vol. 3, no. 2,
pp. 99--114, 1991.

\bibitem{schneier} B. Schneier and J. Kelsey, ``Secure Audit Logs to
Support Computer Forensics,'' \emph{ACM Transactions on Information
and System Security}, vol. 2, no. 2, pp. 159--176, 1999.

\bibitem{rfc6962} B. Laurie, A. Langley, and E. Kasper, ``Certificate
Transparency,'' RFC 6962, IETF, 2013.

\bibitem{prov} W3C, ``PROV-O: The PROV Ontology,'' W3C Recommendation,
2013.

\bibitem{oais} CCSDS, ``Reference Model for an Open Archival
Information System (OAIS),'' ISO 14721:2012.
\end{thebibliography}

\end{document}
